% !TeX root = ../main.tex
\section{Conclusion}
\label{sec:conclusion}

Each part of guitar tone estimation has working methods under easier
conditions. Effects are recognised and their settings regressed on synthetic,
isolated guitar; amplifiers are captured from paired audio to near
transparency and, less accurately, from unpaired audio; tone-transfer models
reproduce whole chains as audio; and inference-time search drives arbitrary,
non-differentiable chains. What is missing is the combination that the
commercial case needs: an interpretable chain, estimated from a separated stem
without the dry signal, and a benchmark that measures it. Products that claim
it publish no method or accuracy, and the closest documented design is a
patent. The evidence points to a hybrid for an effects application: render
the application's own chain as the forward model, predict an initial setting
with a network trained on separated renders of that chain, refine it by
gradient-free search against a pseudo-dry signal, and finish with a linear
match. The transcribed notes, which a tablature system already produces, make
two parts of this easier than in the literature: an aligned resynthesised dry
signal, and the timing needed to measure delay and reverberation tails.
